\subsection{Non-smooth Optimization}
% \paragraph{Convergence Criterion}

Suppose $F$ is differentiable. $\vecx$ is an $\epsilon$-stationary point of $F$ if $\|\nabla F(\vecx)\|\le \epsilon$. This definition is the standard criterion for smooth non-convex optimization.
For non-smooth non-convex optimization, the standard criterion is the following: $\vecx$ is an $(\delta,\epsilon)$-Goldstein stationary point of $F$ if there exists $S\subset \RR^d$ and $P\in\calP(S)$ such that $\vecy\sim P$ satisfies $\E[\vecy]=\vecx$, $\|\vecy-\vecx\|\le\delta$ almost surely, and $\|\E[\nabla F(\vecy)]\|\le\epsilon$.\footnote{The original definition of $(\delta,\epsilon)$ Goldstein stationary point proposed by \cite{zhang2020complexity} does not require the condition $\E[\vecy]=\vecx$. However, as shown in \cite{cutkosky2023optimal}, this condition allows us to reduce a Goldstein stationary point to an $\epsilon$-stationary point when the loss is second-order smooth. Hence we also keep this condition.} Next, we formally define $(c,\epsilon)$-stationary point, our proposed new criterion for non-smooth optimization.

\begin{definition}
\label{def:Exp-O2NC:stationary-point}
Suppose $F:\RR^d\to\RR$ is differentiable, $\vecx$ is a $(c,\epsilon)$-stationary point of $F$ if $\|\nabla F(\vecx)\|_c \le \epsilon$, where
\begin{equation*}
    \|\nabla F(\vecx)\|_c = \inf_{\substack{S\subset \RR^d \\ \vecy\sim P\in\calP(S) \\ \E[\vecy]=\vecx}} \| \E[\nabla F(\vecy)]\| + c \cdot \Ex\|\vecy-\vecx\|^2.
\end{equation*}
\end{definition}

In other words, if $\vecx$ is a $(c,\epsilon)$-stationary point, then there exists $S\subset \RR^d, P\in\calP(S)$ such that $\vecy\sim P$ satisfies $\E[\vecy]=\vecx$, $\E\|\vecy-\vecx\|^2 \le \epsilon/c$, and $\|\E[\nabla F(\vecy)]\|\le \epsilon$. To see how this definition is related to the previous $(\epsilon,\delta)$-Goldstein stationary point definition, consider the case when $c=\epsilon/\delta^2$. Then this new definition of $(c,\epsilon)$-stationary point is almost identical to $(\delta,\epsilon)$-Goldstein stationary point, except that it relaxes the constraint from $\|\vecy-\vecx\|\le\delta$ to $\E\|\vecy-\vecx\|^2\le\delta^2$. 

% ELABORATE why this is good definition

To further motivate this definition, we demonstrate that $(c,\epsilon)$-stationary point retains desirable properties of Goldstein stationary points. Firstly, the following result shows that, similar to Goldstein stationary points, $(c,\epsilon)$-stationary points can also be reduced to first-order stationary points with proper choices of $c$ when the objective is smooth or second-order smooth. 

\begin{restatable}{lemma}{CriterionReduction}
\label{lem:Exp-O2NC:criterion-reduction}
% Assume $\|\nabla F(\vecx)\|_c\le \epsilon$ for any $c=\epsilon\delta^{-2}$ and $\delta>0$. Then $\|\nabla F(\vecx)\| \le \epsilon$ if $F$ is $H$-smooth and $\delta=\eps/2H$ or if $F$ is $\rho$-second-order-smooth and $\delta=\sqrt{2\epsilon}/\sqrt{\rho}$.
Suppose $F$ is $H$-smooth. 
If $\|\nabla F(\vecx)\|_c\le\epsilon$ where $c=H^2\epsilon^{-1}$, then $\|\nabla F(\vecx)\|\le 2\epsilon$. \\
% If $\|\nabla F(\vecx)\|_c\le\epsilon$ where $c=\epsilon\delta^{-2}, \delta=\epsilon/2H$, then $\|\nabla F(\vecx)\|\le 2\epsilon$. \\
Suppose $F$ is $\rho$-second-order-smooth. 
If $\|\nabla F(\vecx)\|_c\le\epsilon$ where $c=\rho/2$, then $\|\nabla F(\vecx)\|\le 2\epsilon$.
% If $\|\nabla F(\vecx)\|_c\le\epsilon$ where $c=\epsilon\delta^{-2}, \delta=\sqrt{2\epsilon}/\sqrt{\rho}$, then $\|\nabla F(\vecx)\|\le 2\epsilon$.
\end{restatable}

As an immediate implication, suppose an algorithm achieves $O(c^{1/2}\epsilon^{-7/2})$ rate for finding a $(c,\epsilon)$-stationary point. Then Lemma \ref{lem:Exp-O2NC:criterion-reduction} implies that, with $c=O(\epsilon^{-1})$, the algorithm automatically achieves the optimal rate of $O(\epsilon^{-4})$ for smooth objectives \cite{arjevani2019lower}. Similarly, with $c=O(1)$, it achieves the optimal rate of $O(\epsilon^{-7/2})$ for second-order smooth objectives \cite{arjevani2020second}. 

Secondly, we show in the following lemma that $(c,\epsilon)$-stationary points can also be reduced to Goldstein stationary points when the objective is Lipschitz.

\begin{restatable}{lemma}{GoldsteinReduction}
\label{lem:Exp-O2NC:goldstein-reduction}
Suppose $F$ is $G$-Lipschitz. For any $c,\epsilon,\delta>0$, a $(c,\epsilon)$-stationary point is also a $(\delta,\epsilon')$-Goldstein stationary point where $\epsilon' = (1+\frac{2G}{c\delta^2})\epsilon$.
\end{restatable}

